\chapter{Evaluation der Methode}
\label{ch:EvaluationDerMethode}
\chaptermark{Evaluation}

\section{Zeitliche Datentrennung}

Alle Teilmengen werden chronologisch statt zufällig gebildet. Imputierte
Zielwerte werden in keiner Teilmenge als Lern- oder Bewertungsbeispiel gezählt.
Tabelle~\ref{tab:forecaster_splits} zeigt die Grenzen nach der Bildung der
Zeitfenster.

\begin{table}[htb]
  \centering
  \caption{Chronologische Teilmengen des Prognosemodells; Grenzen in UTC.}
  \label{tab:forecaster_splits}
  \begin{tabular}{l@{\hspace{1em}}l@{\hspace{1em}}l@{\hspace{1em}}r}
    \hline
    Teilmenge & Beginn & Ende (exklusiv) & gemessene Ziele \\
    \hline
    Training & erster Zielwert & 08.07.2025 & 4.083 \\
    Validierung & 08.07.2025 & 09.07.2025 & 318 \\
    Kalibrierung & 09.07.2025 & 10.07.2025 & 368 \\
    Test & 10.07.2025 & Reihenende & 1.346 \\
    \hline
  \end{tabular}
\end{table}

Skalierungswerte und Modellparameter verwenden nur Training beziehungsweise
Training plus Validierung. Der Kalibrierungsabschnitt dient ausschließlich zur
Bestimmung der Alarmgrenzen. Da der Testzeitraum in früheren
Entwicklungsschritten bereits betrachtet wurde, dient er als zeitlich
abgetrennter Rückblickstest und nicht als völlig unangetasteter Zukunftstest.

\section{Vergleich mit dem Basismodell}

Als Basismodell wird der erste einfache LSTM-Aufbau des Projekts bezeichnet:
\[
 \operatorname{Input}(24,2)\rightarrow\operatorname{LSTM}(64)
 \rightarrow\operatorname{Dense}(2).
\]
Sein damaliger Trainingslauf ist nicht unmittelbar mit dem aktuellen Stand
vergleichbar. Die erste Fensterbildung arbeitete auf 17.578 Rohzeilen und
erzeugte im Test 3.406 Ausgaben für nur 1.346 unterschiedliche Sensorstunden,
weil Mehrfachmessungen noch nicht zusammengefasst waren. Modell,
Skalierungswerte und Schwellen wurden außerdem nicht gemeinsam gespeichert.

Für den fairen Vergleich wird deshalb der Aufbau des Basismodells auf die
heutige korrigierte Datenaufbereitung übertragen. Beide Netze erhalten
dieselben zusammenhängenden 24-Stunden-Fenster, Zielwerte, Zeitgrenzen,
\(\log(1+x)\)-Standardisierung, den Zufallsstartwert 42, Batchgröße 128 und 54
Epochen. Das Basismodell sieht die beiden Messwerte; das aktuelle LSTM erhält
zusätzlich die Beobachtungsmaske. Jedes Netz bestimmt aus demselben Abschnitt
eine eigene 99-\%-Alarmgrenze für jede Messgröße.

Damit werden die beiden vollständigen Modelle verglichen. Der Versuch zeigt
nicht, welchen Anteil eine einzelne Schicht am Unterschied hat.

\section{Kennzahlen für Prognose und Modellgröße}

Für beide LSTM-Modelle sowie \(t-24\), Sensor-Stundenmittel und den Median der
vorhergehenden 24 Stunden werden MAE, RMSE und MASE berechnet. Der MASE-Nenner
ist der mittlere absolute Fehler einer saisonalen \(t-24\)-Prognose auf
gemessenen Trainingszielen \cite{hyndman2006accuracy}:
\begin{equation}
 \operatorname{MASE}=
 \frac{H^{-1}\sum_{h=1}^{H}|y_{T+h}-\hat y_{T+h}|}
 {(T-24)^{-1}\sum_{t=25}^{T}|y_t-y_{t-24}|}.
\end{equation}
MAE und RMSE messen die durchschnittliche beziehungsweise quadratisch stärker
gewichtete Abweichung. Beim MASE bedeutet ein Wert unter eins, dass das
Verfahren besser als die saisonale Trainingsreferenz ist. Zusätzlich zeigt die
Parameterzahl die Größe des Netzes.

\section{Prüfung mit künstlichen Anomalien}

Vollständige reale Kennzeichnungen liegen nicht vor. Deshalb werden bekannte
Messwerte gezielt verändert; diese künstlichen Anomalien ersetzen aber keine
fachliche Prüfung \cite{schmidl2022evaluation,goswami2023unsupervised}.
Untersucht werden kurze Ausschläge nach oben und unten (Spikes und Dips),
sprunghafte Niveauänderungen (Level-Shifts), schleichende Verschiebungen
(Drifts) und festgefrorene Werte (Stuck-at). Die Ereignisse werden unabhängig
vom Modell aus zusammenhängenden gemessenen Abschnitten gewählt. Drei
festgelegte Zufallsstartwerte sorgen bei beiden Modellen für dieselbe Auswahl.

Der direkte Modellvergleich verwendet die starken künstlichen Veränderungen
mit Stärke 2{,}0; für Stuck-at gibt es nur eine Variante. Jede veränderte Reihe
wird mit derselben unveränderten Reihe verglichen, damit bereits vorhandene
Alarme nicht als Treffer zählen. Precision bezeichnet den Anteil richtiger
Alarme, Recall den Anteil erkannter Anomalien und F1 fasst beide Werte
zusammen. Die PR-AUC bewertet die Trennung über verschiedene Alarmgrenzen.
Zusätzlich wird angegeben, bei wie vielen zusammenhängenden Ereignissen
mindestens eine Stunde erkannt wurde. Weitere Stärken sowie ROC-AUC und die
zeitliche Überlappung werden nur in den ausführlichen Ergebnisdateien geführt.

\section{Evaluation der Imputation}

Da die wahren Werte realer Lücken unbekannt sind, werden ursprünglich gemessene
interne Trainingssegmente von einer, zwei, vier und acht Stunden künstlich
entfernt. Vor und nach jeder Testlücke bleibt eine echte Messung erhalten.
Dieselbe Datenaufbereitung rekonstruiert die verdeckten Werte. Drei
Wiederholungen mit festgelegter Zufallsauswahl und jeweils zwölf Ereignissen je
Länge ergeben 540 verdeckte Sensorstunden. Berichtet werden
MAE und RMSE gegen die gesicherten Originalwerte.

\section{Rauschrobustheit und reale Fallprüfung}

Auf die skalierten Eingabewerte wird zufälliges Rauschen mit Normalverteilung und
Standardabweichung 0{,}1 beziehungsweise 0{,}25 gegeben. Der resultierende MAE
wird mit dem Grundlauf verglichen. Die elf Ereignisse eines früheren
Prognoselaufs bleiben als manuelle Fallstudie erhalten. Sieben wurden als
plausibel normal, eines als extern gestütztes Rosengarten-Ereignis und drei als
offen eingeordnet. Da diese Stunden durch ein früheres Verfahren vorausgewählt
wurden, ersetzen sie keine vollständige Prüfung aller Teststunden.

\section{Reproduzierbarkeit}

Die vollständige Auswertung wird in der Projektumgebung \texttt{.venv} in
dieser Reihenfolge ausgeführt: Datenaufbereitung, Imputationsprüfung, Training,
Auswertung und Prüfung mit künstlichen Anomalien. Für den Modellvergleich wird
das oben beschriebene Basismodell einmal unter denselben Bedingungen wie das
aktuelle LSTM trainiert. Ergebnisse aus zwischenzeitlichen Suchläufen werden
nicht berichtet.
